\begin{table}
\caption{Prediction targets of the curated analysis dataset}
\label{tab:table1_dataset}
\tablefont
\begin{tabularx}{\textwidth}{@{}>{\raggedright\arraybackslash}Xlrrlr@{}}
\toprule
\textbf{Target} & \textbf{Unit} & \textbf{Records} & \textbf{Formulas} & \textbf{Train / validation / test} & \textbf{Matched records} \\
\midrule
Magnetic ordering (FM/AFM/FiM/NM) & class & 79,125 & 70,293 & 55,249/11,974/11,902 & 6,061 \\
Curie temperature & K & 29,740 & 14,145 & 20,508/4,691/4,541 & 4,602 \\
DFT total magnetic moment, per cell ($m_\mathrm{cell}$) & $\mu_\mathrm{B}$/cell & 87,513 & 77,402 & 61,254/13,130/13,129 & 10,852 \\
Mass magnetization (secondary) & emu/g & 4,317 & 3,218 & 3,021/648/648 & -- \\
Volume magnetization (secondary) & A/m & 779 & 454 & 535/124/120 & -- \\
Anisotropy constant $K_1$ (secondary) & J/m$^3$ & 1,847 & 1,182 & 1,296/280/271 & -- \\
\bottomrule
\end{tabularx}
\par\vspace{4pt}\raggedright\noindent Analysis dataset v1.4 (frozen); formula-grouped 70/15/15 splits. Matched records = materials with complete Tier 3 descriptors, used at every tier. Ordering classes: FM 34,098, NM 31,673, FiM 7,161, AFM 6,193. The DFT moment is the Materials Project/JARVIS total moment of the computed cell at 0 K (not an experimental saturation magnetization); the moment per formula unit ($m_\mathrm{fu}$ = $m_\mathrm{cell}\times$ atoms per formula unit / atoms in the cell, 10,821 matched rows) is analysed alongside. Tier 1 = 65 composition descriptors; Tier 2 = Tier 1 + 8 oxidation-state descriptors; Tier 3 = Tier 2 + 24 non-magnetic DFT descriptors (97).
\end{table}
